%% Atomic data in EXHALE compared with MoCHII, and the 2026-07-17 rate update.
%% Author: Kwang-Il Seon
\documentclass[english,a4paper]{my_memo}
\synctex=1
\usepackage{graphicx}
\usepackage{booktabs}
\usepackage{amsmath,amssymb}
\usepackage{array}
\usepackage{xcolor}
\usepackage{babel}
\setlength{\emergencystretch}{3em}
% top-align a float that ends up alone on the last page
\makeatletter
\setlength{\@fptop}{0pt}
\makeatother

\newcommand{\acm}{\ensuremath{\,\mathrm{cm^3\,s^{-1}}}}
% zero-width breakpoint, so long typewriter paths can wrap at separators
\newcommand{\bk}{\discretionary{}{}{}}
\newcommand{\code}[1]{\texttt{#1}}

\begin{document}

\title{Atomic Data in EXHALE versus MoCHII:\\
Comparison and the 2026-07-17 Rate Update}
\author{Kwang-Il Seon}
\date{2026-07-17 (updated 2026-08-19)}
\maketitle

\tableofcontents
\bigskip

%=============================================================================
\section{Purpose and scope}
%=============================================================================

This memo has two parts. Part~1 compares the atomic data that EXHALE and
MoCHII (a separate code, not part of this repository) use for the microphysical
processes the two codes have in common: photoionization cross sections,
recombination, collisional ionization, charge exchange, cooling, and
heating. The conclusion is that, where the two codes overlap, they draw on
the same lineage of published data; the differences are in coverage
(MoCHII spans more elements; EXHALE adds the metastable He\,I\,$2^3$S
network and planet-atmosphere species) and in a few modeling choices.

Part~2 documents the rate update applied on 2026-07-17. Before that update
EXHALE inherited its H/He recombination and collisional-ionization fits
from ATES, and several of those fits carried no source annotation. The
update makes the EXHALE default match the Badnell/Mao recombination and
Voronov collisional ionization that MoCHII uses, introduces the
\code{legacy\_hhe\_rates} switch to reproduce the old fits, replaces the
free-free Gaunt factor with the van Hoof et al.\ (2014) thermal average,
and corrects a physical error in the free-free charge weighting.

%=============================================================================
\section{Part 1: Process-by-process comparison}
%=============================================================================

Table~\ref{tab:overlap} summarizes the overlapping processes. The
subsections below add the details that do not fit in the table.

\begin{table}[htbp]
\centering
\footnotesize
\renewcommand{\arraystretch}{1.25}
\caption{Atomic data for the processes EXHALE and MoCHII have in common.
Both codes descend from the same published sources for each process; the
right two columns note where they diverge. EXHALE rows marked
``[updated]'' reflect the 2026-07-17 default (Part~2).}
\label{tab:overlap}
\begin{tabular}{@{}p{0.14\textwidth}p{0.40\textwidth}p{0.38\textwidth}@{}}
\toprule
\textbf{Process} & \textbf{EXHALE} & \textbf{MoCHII} \\
\midrule
Photoionization cross sections &
Verner et al.\ 1996 (VFKY96) with Verner \& Yakovlev 1995 fallback; H\,I
and He\,II from the hydrogenic analytic form. Adds He\,I\,$2^3$S (VFKY96
wing bridge) and H$_2$ (Yan, Sadeghpour \& Dalgarno 1998). Tuned for
planet atmospheres (Na, K). &
Verner et al.\ 1996 with Verner \& Yakovlev 1995 fallback; H\,I and
He\,II hydrogenic. Wider element set (Ne, Cl, Ar). \\

Recombination (H, He) &
[updated] Badnell (2023) radiative $\alpha_A$ minus Mao \& Kaastra (2016)
$\alpha_1$, giving case~B; He\,II includes Badnell dielectronic
recombination. Identical coefficients to MoCHII. &
Badnell (2023) radiative $\alpha_A$ and Mao \& Kaastra (2016) $\alpha_1$
for case~B, with Badnell dielectronic recombination for He\,II. \\

Recombination (metals) &
Badnell radiative/dielectronic recombination (Badnell 2006 and adf48
data), Ca\,I from Shull \& Van Steenberg 1982, Fe from the Huang et al.\
2023 fit (Badnell covers neither the K-like nor the Mn-like/Cr-like
sequences). &
Badnell radiative/dielectronic recombination
(\code{badnell\_rr/dr.dat}) with a CHIANTI fallback. \\

Collisional ionization (H, He) &
[updated] Voronov 1997. &
Voronov 1997 (unified fits), with Dere 2007 as an option. \\

Collisional ionization (metals) &
Voronov 1997. &
Voronov 1997. \\

Charge exchange &
Huang et al.\ 2023 (ApJ 951, 123) Table~4, based on Kingdon \& Ferland
1996. 63 reactions. On by default: the 23 metal$+$H reactions of Group~A,
and the He$\leftrightarrow$H pair of Group~B under its own key
\code{He\_H\_}\bk{}\code{charge\_}\bk{}\code{exchange} (default \code{True},
independent of \code{cx\_full}). \code{cx\_full} adds the
metal$+$He and metal$+$metal reactions of Groups~C--D. &
Huang et al.\ 2023 Table~4 (Kingdon \& Ferland 1996), metals only. \\

Recombination cooling &
Hui \& Gnedin 1997 case-B cooling fits. &
Hui \& Gnedin 1997 lineage, with a case-A/case-B choice. \\

Free-free cooling &
[updated] Constant $1.426\times10^{-27}$; thermally-averaged Gaunt factor
from van Hoof et al.\ 2014 (161-point table). &
Constant $1.426\times10^{-27}$; Gaunt from Hummer 1988 / van Hoof et al.\
2014. \\

Metal line cooling &
CHIANTI v11.0.2 fit $\Lambda(T)=T^{-1/2}\sum_i A_i e^{-T_i/T}$
(\code{cooling\_data/fit\_cno\_formulas.py}), extended with
density-dependent treatments (see \S\ref{sec:unique}). &
CHIANTI v11.0.2 fit of the same form (\code{chianti\_cooling.py}), in the
low-density limit; emission lines handled by a separate n-level solver. \\

Heating (photoelectric) &
[updated] Shull \& van Steenberg (1985) secondary ionization in the main
loop: a photoelectron with $E_0>40$~eV deposits only $f_{\rm heat}(x)$ of its
excess as heat and drives H\,I / He\,I secondary ionizations (default on,
\code{use\_sec\_ion}). Below 40~eV the photoelectron thermalizes fully. &
Full photoelectron excess energy thermalized; no secondary ionization. \\
\bottomrule
\end{tabular}
\end{table}

\subsection{Photoionization cross sections}
Both codes use Verner et al.\ (1996) fits with the Verner \& Yakovlev
(1995) inner-shell data as a fallback, and both take H\,I and He\,II from
the hydrogenic analytic form. EXHALE adds two cross sections that MoCHII
does not carry: the metastable He\,I\,$2^3$S level (bridged onto the
VFKY96 wing) and molecular hydrogen (Yan, Sadeghpour \& Dalgarno 1998).
MoCHII covers a broader element list (including Ne, Cl, and Ar); EXHALE
instead carries the alkali species (Na, K) relevant to planet
atmospheres.

\subsection{Recombination}
Before 2026-07-17 the EXHALE H/He case-B coefficients were unattributed
fits inherited from ATES. A numerical check showed them to be exactly the
Hui \& Gnedin (1997) case-B fits. MoCHII builds its default case-B
coefficients from the Badnell (2023) radiative $\alpha_A$ and the Mao \&
Kaastra (2016) $\alpha_1$ subtraction, with Badnell dielectronic
recombination added for He\,II. The 2026-07-17 update (Part~2) adopts the
same construction as the EXHALE default, so the two codes now agree on the
H/He recombination coefficients. For metals, both codes already used
Badnell radiative and dielectronic recombination, with small differences
in the fallback data for a few ions.

\subsection{Collisional ionization}
Before the update the EXHALE H\,I and He\,I collisional-ionization rates
came from Abel et al.\ (1997), which reproduces the Janev et al.\ (1987)
fits, and He\,II from a Hui \& Gnedin (1997) fit; none of these carried a
source annotation. MoCHII uses the unified Voronov (1997) fits throughout,
with Dere (2007) as an option. The 2026-07-17 update makes Voronov (1997)
the EXHALE default for H and He as well; the metal rates were already
Voronov (1997) in both codes.

\subsection{Charge exchange}
Both codes transcribe Table~4 of Huang et al.\ (2023, ApJ 951, 123), which
rests on Kingdon \& Ferland (1996). EXHALE carries 63 of its 65 rows (23 in
Group~A, 2 in B, 6 in C, 32 in D; the two missing ones are in the
metal--metal set). Two switches select them. The 23 metal$+$H reactions of
Group~A are on by default and \code{cx\_full} adds the metal$+$He (C) and
metal$+$metal (D) reactions. The He$\leftrightarrow$H pair (B1, B2) is separate:
it is not in the generic reaction list at all but is applied by dedicated
routines in every ionization system containing He, gated by its own input key
\code{He\_H\_charge\_exchange} (default \code{True}) and independent of
\code{cx\_full}. MoCHII carries the metal reactions only.

The printed Table~4 exchanges the reactant labels of its two oxygen rows:
the $\exp(-227/T)$ Boltzmann factor appears on $\mathrm{O^+ + H^0}$, but
$\mathrm{IP(O\,\textsc{i})} > \mathrm{IP(H\,\textsc{i})}$ makes
$\mathrm{O^0 + H^+}$ the endothermic direction that must carry it. EXHALE
assigns the two rate coefficients to the physically correct rows; see
\code{HUANG2023\_TABLE4\_OXYGEN\_ERRATUM.md} for the detailed-balance and
Cloudy cross-checks. Any code transcribing the table as printed inherits
the error.

\subsection{Cooling}
Recombination cooling in both codes follows the Hui \& Gnedin (1997)
case-B family (MoCHII offers a case-A/case-B choice). Both use the same
free-free constant, $1.426\times10^{-27}$; the Gaunt factor is discussed
in \S\ref{sec:gaunt}. Metal line cooling in both codes is a CHIANTI
v11.0.2 fit of the form $\Lambda(T)=T^{-1/2}\sum_i A_i e^{-T_i/T}$
(EXHALE: \code{cooling\_data/fit\_cno\_formulas.py}; MoCHII:
\code{chianti\_cooling.py}). The MoCHII fit is the low-density limit, with
the emission lines themselves computed by a separate n-level solver.

\subsection{Heating}
EXHALE now includes Shull \& van Steenberg (1985) secondary ionization by
fast photoelectrons in the main loop (default on, \code{use\_sec\_ion}; see
Part~2). MoCHII thermalizes the full excess energy of the photoelectron and
does not add secondary ionization, which is appropriate in the nearly fully
ionized H\,II region it targets. The \code{wind\_ae} initial-condition
generator separately uses the same SvS85 partition together with the Dere
(2007) coefficients.

\subsection{Processes carried by only one code}\label{sec:unique}
EXHALE carries several networks that MoCHII does not: the metastable
He\,I\,$2^3$S chemistry (Oklop\v{c}i\'c \& Hirata 2018; Bray 2000; Penning
ionization from Taylor et al.\ 2025), the H\,($n{=}2$) 2s/2p populations
(Christie, Arras \& Li 2013; Draine 2011), Lyman-$\alpha$ radiative
transfer (Neufeld 1990), and the H$_2$/H$_3^+$ molecular network (Miller
et al.\ 2013). Its cooling also adds density-dependent treatments absent
from the MoCHII low-density fits: a two-dimensional $(T,n_e)$ statistical
equilibrium table for Fe\,II, the ground terms of C\,I, C\,II, N\,II and O\,I
solved in statistical equilibrium (the origin of the [C\,II] 158\,$\mu$m and
[O\,I] 63\,$\mu$m saturation), the Mg\,II h\&k, Ca\,II H\&K, and Na\,I~D two-level
terms, and a Fe\,I table from NIST.

MoCHII in turn carries an n-level statistical-equilibrium emission-line
solver (CHIANTI v11), Storey \& Hummer (1995) recombination lines, Porter
et al.\ (2012, 2013) He\,I case-B emissivities, the nebular continuum
(free-bound plus two-photon, from Nussbaumer \& Schmutz 1984 / Almog \&
Netzer 1989), and dust photoelectric heating (Bakes \& Tielens 1994).

One quantity crosses between them. MoCHII transports the He\,I $2^1$S
two-photon continuum as sampled packets and so needs the full Drake, Victor
\& Dalgarno (1969) shape; EXHALE never forms a spectrum, but its He\,II
cascade coupling (\texttt{he\_rec\_coupling}) needs two moments of that same
shape above the H\,I edge: how many ionizing photons a $2^1$S decay yields
and how much energy each of them deposits. Both had been round numbers, 0.56
and 3.0\,eV. Integrating the shape gives 0.5564 and 2.512\,eV: the photon
count was right, but the deposited energy was 20\% high, being close to the
3.511\,eV a uniform in-band distribution would give rather than to the peaked
one the pair actually has. MoCHII found the same defect in its own sampler,
which drew the in-band photons flat. EXHALE now carries the two integrals as
\texttt{f\_2q\_HeI} and \texttt{Ee\_2q\_HeI}.

%=============================================================================
\section{Part 2: The 2026-07-17 rate update}
%=============================================================================

The changes below are implemented in the radiation modules
\code{Cool\_coeff.f90} and \code{util\_ion\_eq.f90}, in the temperature
solver \code{T\_equation.f90}, and in the input parsing under
\code{src/modules/files\_IO/}. Items~1, 2, and 3 are gated by the new
\code{legacy\_hhe\_rates} switch; items~4 and 5 apply always, independent
of the switch.

\subsection{H/He recombination default}
The default H/He case-B coefficients are now
\begin{equation}
\alpha_B = \alpha_A^{\rm Badnell} - \alpha_1^{\rm Mao},
\end{equation}
where $\alpha_A^{\rm Badnell}$ is the Badnell radiative total (with the
three-term Badnell dielectronic recombination added for He\,II) and
$\alpha_1^{\rm Mao}$ is the Mao \& Kaastra (2016, A\&A 587, A84)
ground-state coefficient. The implementation is in
\code{Cool\_coeff.f90}: \code{rr\_badnell}, \code{dr\_HeII\_badnell},
\code{rr\_mao}, and the assembled \code{alphaB\_HII\_new} /
\code{alphaB\_HeII\_new} / \code{alphaB\_HeIII\_new}. The coefficients match
those in the MoCHII \code{recomb\_mod.f90}.

\subsection{H/He collisional ionization default}
The default H/He collisional-ionization rates are now the Voronov (1997,
ADNDT 65, 1) fits, implemented through \code{voronov\_ci}
(\code{Cool\_coeff.f90}). The adopted parameters are listed in
Table~\ref{tab:voronov}.

\begin{table}[htbp]
\centering
\small
\caption{Voronov (1997) collisional-ionization parameters adopted as the
EXHALE default.}
\label{tab:voronov}
\begin{tabular}{@{}lccccc@{}}
\toprule
Ion & $\Delta E$ [eV] & $P$ & $A$ & $X$ & $K$ \\
\midrule
H\,I   & 13.6 & 0 & $2.91\times10^{-8}$ & 0.232 & 0.39 \\
He\,I  & 24.6 & 0 & $1.75\times10^{-8}$ & 0.180 & 0.35 \\
He\,II & 54.4 & 1 & $2.05\times10^{-9}$ & 0.265 & 0.25 \\
\bottomrule
\end{tabular}
\end{table}

\subsection{The \code{legacy\_hhe\_rates} switch}
A new input flag \code{legacy\_hhe\_rates} (default \code{False}) restores
the previous fits: Hui \& Gnedin (1997) recombination and the
Abel et al.\ (1997) / Hui \& Gnedin (1997) collisional ionization. It is
declared in \code{parameters.f90}, defaulted and parsed in
\code{input\_read.f90}, and echoed to \code{parse\_dump.txt} by
\code{write\_setup\_report.f90}: grep the flag name in each; the line numbers
this paragraph used to quote had gone stale. It
is registered as row K14b in \code{docs/input\_schema.md}. Setting it to
\code{True} reverts items~1 and 2 exactly; the free-free changes (items~4
and 5) are not affected.

\subsection{Free-free Gaunt factor}\label{sec:gaunt}
The free-free Gaunt factor is now the thermally-averaged non-relativistic
$\langle g_{\rm ff}\rangle(\gamma^2)$ of van Hoof et al.\ (2014, MNRAS
444, 420), with $\gamma^2 = Z_{\rm ion}^2\,\mathrm{Ry}/kT$. The frequency
average
\begin{equation}
\langle g_{\rm ff}\rangle(\gamma^2)
  = \frac{\int e^{-u}\,g_{\rm ff}(\gamma^2,u)\,du}{\int e^{-u}\,du}
\end{equation}
is tabulated at 161 points in $\log_{10}\gamma^2 = -6$ to $10$ in steps of
0.1~dex (generated by \code{cooling\_data/\bk gauntff\_\bk thermal\_\bk avg.py}; the
table is \code{gff\_avg} in \code{Cool\_coeff.f90}). This
replaces the previous unattributed two-branch logarithmic fit and applies
always, regardless of \code{legacy\_hhe\_rates}. The tabulated value at
$T=10^4$\,K and $Z=1$ is $\langle g_{\rm ff}\rangle = 1.266$.

\subsection{Free-free charge-weighting correction}
The previous code weighted He\,II (net charge $+1$) by $Z^2=4$ and
evaluated the metal Gaunt factor at the nuclear atomic number ($Z=8$ for
O, $Z=26$ for Fe). This is physically wrong: free-free emission scales
with the ion \emph{net} charge $Z_{\rm ion}$, not the nuclear number. The
correction sets $Z_{\rm ion}=1$ for H\,II, He\,II, and singly-ionized
metals, and $Z_{\rm ion}=2$ for He\,III and doubly-ionized metals
(\code{util\_ion\_eq.f90} \code{eval\_cool}, lines~376--396;
\code{T\_equation.f90}, lines~105--121). The most visible consequence is
that the free-free cooling of a He\,II-dominated layer drops by a factor
of four. Like the Gaunt-table change, this correction applies always.

\subsection{Benchmark values}
Table~\ref{tab:bench} lists the new default case-B recombination and
collisional-ionization rates at $T=10^4,\,10^5,\,10^6$\,K, with the legacy
values at $10^4$\,K where relevant. The Fortran implementation reproduces
an independent Python recomputation to all printed digits, and the legacy
branch reproduces the pre-update values exactly.

\begin{table}[!ht]
\centering
\small
\caption{New default rates ($\alpha_B$ and collisional ionization $ci$;
units \acm{}), with legacy values at $10^4$\,K for comparison. The legacy
$\alpha_B(\mathrm{H\,II})$ at $10^4$\,K is the Hui \& Gnedin (1997) fit,
$2.592\times10^{-13}$.}
\label{tab:bench}
\begin{tabular}{@{}lcccc@{}}
\toprule
Coefficient & $T=10^4$\,K & $T=10^5$\,K & $T=10^6$\,K & legacy ($10^4$\,K)\\
\midrule
$\alpha_B(\mathrm{H\,II})$  & $2.606\times10^{-13}$ & $3.070\times10^{-14}$ & $2.163\times10^{-15}$ & $2.592\times10^{-13}$ \\
$\alpha_B(\mathrm{He\,II})$ & $2.807\times10^{-13}$ & $5.006\times10^{-13}$ & $1.015\times10^{-12}$ & $2.616\times10^{-13}$ \\
$\alpha_B(\mathrm{He\,III})$& $1.536\times10^{-12}$ & $2.370\times10^{-13}$ & $2.244\times10^{-14}$ & - \\
$ci(\mathrm{H\,I})$         & $7.448\times10^{-16}$ & $3.963\times10^{-9}$  & $3.103\times10^{-8}$  & $7.247\times10^{-16}$ \\
\bottomrule
\end{tabular}
\end{table}

The new $\alpha_B(\mathrm{H\,II})$ at $10^4$\,K differs from the legacy
Hui \& Gnedin (1997) value by about $1\%$; the He\,II and He\,III case-B
coefficients differ more, mainly because the new He\,II value carries the
Badnell dielectronic contribution explicitly.

\subsection{Secondary ionization in the main loop}\label{sec:secion}
A photoelectron ejected with kinetic energy $E_0 = e_\nu - E_{\rm th}$ above
about 40~eV does not thermalize immediately: before it does, it collisionally
ionizes further H\,I and He\,I atoms. Shull \& van Steenberg (1985) give
asymptotic fits for the partition of the photoelectron's energy as a function
of the ionized fraction $x$ of the H+He nuclei,
\begin{align}
  f_{\rm heat}(x)    &= 0.9971\,(1 - (1 - x^{0.2663})^{1.3163}), \\
  f_{\rm ion,HI}(x)  &= 0.3908\,(1 - x^{0.4092})^{1.7592}, \\
  f_{\rm ion,HeI}(x) &= 0.0554\,(1 - x^{0.4614})^{1.6660}.
\end{align}
As $x\to1$ the heating fraction approaches unity and both ionization fractions
go to zero; as $x\to0$ the heating fraction vanishes and the ionization
fractions reach their maxima (0.3908 and 0.0554). In neutral gas a large share
of the photoelectron energy therefore goes into secondary ionizations rather
than heat.

A 40~eV threshold gates the partition: for $E_0\le40$~eV the photoelectron
thermalizes fully (no secondary ionization), matching the \code{wind\_ae}
X-ray cutoff. SvS85 is strictly an $E_0\gtrsim100$~eV asymptotic fit, so
applying it down to 40~eV is a deliberate approximation adopted here.

Only H\,I and He\,I receive secondary ionizations, the two channels SvS85
resolves. The SvS85 Ly$\alpha$ excitation channel is assumed to escape as line
radiation and is not put into any rate (a future refinement could couple it to
the Ly$\alpha$ field of the excited-H model). Each absorbing species $s$
(H\,I, He\,I, He\,II, the He\,I triplet when active, H$_2$ when present, and
the photoionizable metal ions)
contributes photoelectrons of energy $E_0 = e_\nu - E_{{\rm th},s}$ that drive
$f_{\rm ion,HI}\,E_0/E_{\rm th,HI}$ secondary H\,I ionizations and
$f_{\rm ion,HeI}\,E_0/E_{\rm th,HeI}$ secondary He\,I ionizations. The He\,I
triplet ($2^3$S, threshold 4.8~eV) previously contributed opacity and its
photoionization rate but no photoheating; its photoelectron energy now enters
the heating and absorbed-energy budgets and the secondary source like every
other absorber.

This differs from how \code{wind\_ae} handles the same physics.
\code{wind\_ae} combines the SvS85 $f_{\rm heat}$ / $f_{\rm exc}$ partition
with the Dere (2007) coefficient tables to split the secondary ionizations
across species, but those Dere tables are tied to the fixed \code{wind\_ae}
spectral grid and cannot be reused in the main loop, which runs on a runtime
SED and a variable energy grid. The main loop therefore applies the
species-resolved SvS85 fits ($f_{\rm ion,HI}$, $f_{\rm ion,HeI}$) directly.

The feature is gated by \code{use\_sec\_ion} (default True). With
\code{use\_sec\_ion = False} the photoelectron thermalizes fully and the
heating and photoionization rates are bit-identical to the legacy path. It is
implemented in the \code{PH\_heat\_H} and \code{PH\_heat\_HHe} routines of
\code{util\_ion\_eq.f90}, wired through the callers in
\code{ionization\_equilibrium.f90} and \code{post\_process\_adv.f90}; the flag
is declared in \code{parameters.f90}, parsed in \code{input\_read.f90}, echoed
by \code{write\_setup\_report.f90}, and documented as key K14c in
\code{input\_schema.md}.

\subsection{Effect on existing results}
Because the default rates change, results computed with the new default
differ from the pre-update golden and gate outputs. Setting
\code{legacy\_hhe\_rates = True} reproduces the old H/He rates and
recovers the previous ionization balance, with two exceptions that apply
in all cases: the free-free charge-weighting correction (\S3.5) and the
van Hoof et al.\ (2014) Gaunt table (\S3.4). Those two changes are
physical corrections and are intentionally not tied to the switch.
Secondary ionization (\S\ref{sec:secion}) is a further default change on top
of the rate update; \code{use\_sec\_ion = False} restores the
full-thermalization heating and photoionization rates bit-identically.

\end{document}
